\begin{proof}[Proof of \cref{obj:lem:continuous-chebyshev-endpoint-bound}]
Put
\[
x_0=\frac{2}{q_{\max}}-1,
\qquad
\lambda_0=x_0+\sqrt{x_0^2-1}.
\]
Since \(0<q_{\max}<1\), one has \(x_0>1\), and hence \(\lambda_0>1\). Choose
\[
C(q_{\max})=2,
\qquad
c(q_{\max})=\frac12(1-\lambda_0^{-1})^2 .
\]
These constants are positive, because \(\lambda_0>1\). % lean: continuous_chebyshev_endpoint_bound
Now let \(\beta\ge1\) and \(0<q\le q_{\max}\), and write
\[
x_q=\frac2q-1,
\qquad
\lambda(q)=x_q+\sqrt{x_q^2-1}.
\]
Then \(q<1\), so \(x_q>1\) and \(\lambda(q)>1\). Moreover the map \(q\mapsto x_q+\sqrt{x_q^2-1}\) is decreasing on this range, so \(q\le q_{\max}\) gives
\[
\lambda_0\le \lambda(q).
\] % lean: endpoint_lambda_mono
Let \(R\) be a real polynomial with \(\deg(R)\le\beta\) and \(|R(x)|\le1\) for \(x\in[-1,1]\). By \cref{obj:lem:chebyshev-exterior-extremal}, applied at the exterior point \(x_q>1\),
\[
|R(x_q)|\le T_\beta(x_q).
\] % lean: chebyshev_exterior_extremal
Also \(|R(-1)|\le1\), and therefore
\[
|R(x_q)-R(-1)|
\le |R(x_q)|+|R(-1)|
\le T_\beta(x_q)+1.
\] % lean: continuous_chebyshev_endpoint_bound
For \(x_q\ge1\), the Chebyshev evaluation formula gives
\[
T_\beta(x_q)
=
\frac{\lambda(q)^\beta+\lambda(q)^{-\beta}}2.
\] % lean: chebyshev_eval_eq_lambda_average
Since \(\lambda(q)\ge1\) and \(\beta\ge1\),
\[
\frac{\lambda(q)^\beta+\lambda(q)^{-\beta}}2+1
\le 2\lambda(q)^\beta.
\] % lean: chebyshev_lambda_average_add_one_le
Combining these inequalities proves
\[
|R(x_q)-R(-1)|\le 2\lambda(q)^\beta
= C(q_{\max})\lambda(q)^\beta.
\]
For the lower bound, the same evaluation formula and the inequalities \(\lambda_0>1\), \(\lambda_0\le\lambda(q)\), and \(\beta\ge1\) give
\[
T_\beta(x_q)-1
=
\frac{\lambda(q)^\beta+\lambda(q)^{-\beta}}2-1
\ge
\frac12(1-\lambda_0^{-1})^2\lambda(q)^\beta
=
c(q_{\max})\lambda(q)^\beta.
\] % lean: chebyshev_lambda_average_sub_one_ge
This proves both assertions. % lean: continuous_chebyshev_endpoint_bound
\end{proof}
